\section{Related Work}

\paragraph{Probabilistic forecast evaluation.}
The Brier score evaluates squared error for probabilistic outcomes \citep{brier1950verification}.
Strictly proper scoring rules reward honest probabilities rather than thresholded classifications
\citep{gneiting2007strictly}. We use Brier score as the primary metric and log loss as a key
secondary metric. Because events are rare, we also report precision--recall area as a descriptive
ranking measure, but it does not replace calibration-sensitive scoring.

\paragraph{Out-of-sample and rolling evaluation.}
Forecast evaluation must preserve temporal order \citep{tashman2000outofsample}. Random train--test
splits would expose a model to future policy regimes and later points in the same event cycle.
Leakage can also arise from feature construction, source timestamps, or post-event explanations
\citep{kapoor2023leakage}. Our expanding-window procedure reconstructs features at each issue time
and evaluates all models on identical timestamps.

\paragraph{Events and organizational attention.}
Event system theory characterizes events by novelty, disruption, and criticality
\citep{morgeson2015event}. The attention-based view argues that organizational action depends on
which issues become salient to decision makers \citep{ocasio1997attention}. These theories motivate
our context variables, but they do not license causal claims. Official incidents can be associated
with both actions and non-actions, and public status records are an incomplete trace of internal
attention.

\paragraph{Drift and reflexivity.}
Changing products and policies create concept drift \citep{gama2014survey}. In addition, publishing
a forecast about a public actor may itself affect discourse or behavior. We therefore preserve model
versions, record policy regimes, and forbid silent replacement of historical predictions. The
resulting study is closer to an audit-preserving longitudinal diagnostic than to a stationary
benchmark.
